\section{Discussion}\label{sec:discuss}

The results support the hypothesis that pseudo-relevance feedback improves
top-10 retrieval effectiveness on Cranfield. MAP@10 increased from 0.3247 to
0.3639, and the paired AP@10 comparison showed substantially more wins than
losses. This suggests that the top retrieved document often contains useful
technical vocabulary that is missing from the original query. In such cases,
feedback expansion helps bridge vocabulary mismatch without requiring external
resources.

The improvement is not uniform. MRR@10 decreased slightly from 0.7759 to 0.7623.
This indicates that feedback can sometimes move the first relevant document
down, especially when the top initial document is only partially relevant or
introduces terms from a narrower subtopic. One failure case was query 77, which
asks about comparing shock-layer theory with experiments in the low Reynolds
number regime. Its AP@10 decreased after feedback, suggesting query drift from
the initially retrieved document. In contrast, queries such as query 20 on Joule
heating in magnetohydrodynamic free convection, query 64 on predicting flutter,
query 78 on the surge line of an axial compressor, and query 91 on transonic
interference effects improved strongly because the feedback document added
technical terms that matched other relevant Cranfield documents.

The baseline VSM limitations remain important. First, VSM depends on lexical
overlap, so it struggles when a query and relevant document use different
technical expressions. Second, it does not understand word sense; the toy
example with \emph{star} shows that the same term can refer to astronomy or
Hollywood. Third, out-of-vocabulary terms receive no vector weight. If a query
contains only terms absent from the document vocabulary, the implemented system
cannot compute a meaningful similarity score and returns the default ordering.
For example, a query containing a misspelled or unseen term such as
``hypersonicx flutterzz'' would have no useful overlap unless spelling
correction or character-level matching were added.

The current approach is therefore best understood as a strong lexical baseline
with a controlled feedback mechanism. It is interpretable, efficient, and
measurably better than the no-feedback baseline for MAP@10, but it is not a full
semantic search system. Methods such as LSA, query spelling correction, or
neural sentence embeddings could address some remaining failures, but they
would require additional validation to avoid replacing one kind of query drift
with another.
